\section{For Ever and Ever}

Every pain we know has an end. The worst illness breaks or kills; the longest grief is softened by years; even a life sentence is measured by a life. It is the last word of Christ's sentence that removes every such measure. The judgment scene closes: ``And these shall go into everlasting punishment: but the just, into life everlasting.''\endnote{Matthew 25:46, Douay--Rheims (Challoner).} On the last Sunday of the year the 1962 Missal reads the Lord's own warrant for that word: ``Heaven and earth shall pass: but my words shall not pass.''\endnote{Matthew 24:35, Douay--Rheims (Challoner), in the Gospel of the Last Sunday after Pentecost in the 1962 Missal, Matthew 24:15--35.} The sentence of the Judge belongs to the words that do not pass.

Augustine saw that the verse answers its own question. The punishment and the life are named in one breath by one Lord with one word, and they stand or fall together:

\begin{quote}
Then what a fond fancy is it to suppose that eternal punishment means long-continued punishment, while eternal life means life without end, since Christ in the very same passage spoke of both in similar terms in one and the same sentence, ``These shall go away into eternal punishment, but the righteous into life eternal!'' If both destinies are ``eternal,'' then we must either understand both as long-continued but at last terminating, or both as endless. For they are correlative,---on the one hand, punishment eternal, on the other hand, life eternal. And to say in one and the same sense, life eternal shall be endless, punishment eternal shall come to an end, is the height of absurdity. Wherefore, as the eternal life of the saints shall be endless, so too the eternal punishment of those who are doomed to it shall have no end.\endnote{Augustine, \emph{City of God} XXI.23.}
\end{quote}

The same Gospel and the same apostles speak the same way in other places. The Lord says of the sin against the Holy Ghost that the one who commits it ``shall never have forgiveness, but shall be guilty of an everlasting sin.''\endnote{Mark 3:29, Douay--Rheims (Challoner).} He says that it is better to enter life maimed ``than having two hands or two feet, to be cast into everlasting fire.''\endnote{Matthew 18:8, Douay--Rheims (Challoner).} He says of his betrayer, ``It were better for him, if that man had not been born'';\endnote{Matthew 26:24, Douay--Rheims (Challoner).} and of no one whose loss would end at last in glory could it be better never to have been born. Saint Paul says that those who know not God and obey not the gospel ``shall suffer eternal punishment in destruction, from the face of the Lord and from the glory of his power.''\endnote{2 Thessalonians 1:9, Douay--Rheims (Challoner).} And the Apocalypse, which closes the Scriptures, sees the devil ``cast into the pool of fire and brimstone, where both the beast and the false prophet shall be tormented day and night for ever and ever,'' and the smoke of the torments of those who worshipped the beast ascending ``for ever and ever: neither have they rest day nor night.''\endnote{Apocalypse 20:9--10 and 14:11, Douay--Rheims (Challoner), whose verse division falls within the sentence at 20:9--10.}

Gregory the Great answered those who said the Lord threatened eternal pain only to frighten men from sin. A threat that God did not mean would make a liar of him, and his promises would be no surer than his threats:

\begin{quote}
Certain it is, and without all doubt most true, that as the good shall have no end of their joys, so the wicked never any release of their torments\ldots{} Seeing, then, true it is, that which he hath promised to his friends: out of all question false it cannot be, that which he hath threatened to his enemies\ldots{} For if he threatened that which he meant not to put into execution: whiles we are desirous to make him merciful, enforced we are likewise (which is horrible to speak) to affirm him to be deceitful.\endnote{Gregory, \emph{Dialogues} IV.44.}
\end{quote}

The Church made that answer her own. When the opinion of a final restoration of the demons and the wicked, associated with the name of Origen and found also in a chapter of Gregory of Nyssa's \emph{Great Catechism},\endnote{Gregory of Nyssa, \emph{The Great Catechism} 26, \emph{Nicene and Post-Nicene Fathers}, 2nd series, vol.~5, Christian Classics Ethereal Library, accessed 16 October 2026.} was condemned at Constantinople in 543, the ninth anathema struck at anyone who ``says or thinks that the punishment of demons and of impious men is only temporary, and will one day have an end.''\endnote{Anathemas against Origen, published in the synod of Constantinople of 543, canon 9, DS 411; English in H. R. Percival, \emph{The Seven Ecumenical Councils}, \emph{Nicene and Post-Nicene Fathers}, 2nd series, vol.~14 (New York, 1900), ``The Anathematisms of the Emperor Justinian against Origen,'' IX.} The Fourth Lateran Council professed the punishment \latin{perpetuam}, and the \emph{Catechism} speaks of separation from God ``for ever.''

\subsection{Why it cannot end}

Why should a sin committed in a moment be punished without end? The Fathers gave two answers, and the Doctors joined them.

The first looks at the will of the sinner. Gregory says that the wicked made an end of sinning only because they made an end of living: ``for willingly they would, had it been in their power, have lived without end, that they might in like manner have sinned without end.''\endnote{Gregory, \emph{Dialogues} IV.44.} The soul that dies in mortal sin has fixed its last end in something other than God. It has chosen, and death seals the choice. Aquinas takes up Gregory's words and explains them: ``A man is said to have sinned in his own eternity, not only as regards continual sinning throughout his whole life, but also because, from the very fact that he fixes his end in sin, he has the will to sin, everlastingly.''\endnote{Aquinas, \emph{Summa theologiae} I-II, q.~87, a.~3, reply to objection 1, citing Gregory, \emph{Dialogues} IV.44.}

The second looks at the wound itself. In the same article Aquinas compares mortal sin to the destruction of an eye:

\begin{quote}
Now disturbance of an order is sometimes reparable, sometimes irreparable: because a defect which destroys the principle is irreparable, whereas if the principle be saved, defects can be repaired by virtue of that principle. For instance, if the principle of sight be destroyed, sight cannot be restored except by Divine power; whereas, if the principle of sight be preserved, while there arise certain impediments to the use of sight, these can be remedied by nature or by art\ldots{} Consequently if a sin destroys the principle of the order whereby man's will is subject to God, the disorder will be such as to be considered in itself, irreparable, although it is possible to repair it by the power of God. Now the principle of this order is the last end, to which man adheres by charity. Therefore whatever sins turn man away from God, so as to destroy charity, considered in themselves, incur a debt of eternal punishment.\endnote{Aquinas, \emph{Summa theologiae} I-II, q.~87, a.~3, corpus.}
\end{quote}

\noindent Mortal sin kills charity, and charity is the very principle by which the soul would return. Such a wound does not heal from within. It can be healed only from without, by the power of God; and in this life God offers that power in the grace of conversion and the sacrament of penance. Death closes the offer. The ancient homily known as the Second Letter of Clement says it with the image of the potter:

\begin{quote}
For as the potter, if he make a vessel, and it be distorted or broken in his hands, fashions it over again; but if he have before this cast it into the furnace of fire, can no longer find any help for it: so let us also, while we are in this world, repent with our whole heart of the evil deeds we have done in the flesh, that we may be saved by the Lord, while we have yet an opportunity of repentance. For after we have gone out of the world, no further power of confessing or repenting will there belong to us.\endnote{Second Letter of Clement 8, trans.\ John Keith, \emph{The Ante-Nicene Fathers}, vol.~9, as revised for New Advent, accessed 16 October 2026.}
\end{quote}

\noindent Gregory Nazianzen describes the judgment that follows, ``after which there is no appeal, no higher court, no defence on the ground of subsequent conduct, no oil obtained from the wise virgins, or from them that sell, for the lamps going out, no repentance of the rich man wasting away in the flame.''\endnote{Gregory Nazianzen, \emph{Oration} 16.9.}

\subsection{The weight of eternity}

Alphonsus gave his twenty-seventh consideration to this alone, and he began it bluntly: ``WERE Hell not eternal, it should not be Hell.'' He did not offer the eternity of hell as a theory: ``This eternity is of faith: it is not an opinion, but a truth attested by God in so many places in the holy Scriptures.''\endnote{Alphonsus, \emph{Preparation for Death}, Consideration XXVII, first point, pp.~328--329.} Then he set the mind to work on it. ``In Hell there is a gate for admission, but none for egress.'' The damned feel not only each moment's pain, but in each moment the whole of what is to come: ``Hence the damned not only suffer the torments of each moment, but in each moment they endure the pain of eternity, saying: What I now suffer, I must suffer for ever.''\endnote{Ibid., second point, pp.~331--332.}

The \emph{Imitation of Christ} measures it against our penances: ``An hour's pain there shall be more grievous than a hundred years here of the bitterest penitence. No quiet shall be there, no comfort for the lost, though here sometimes there is respite from pain, and enjoyment of the solace of friends.''\endnote{Thomas \`a Kempis, \emph{The Imitation of Christ} I.24.4, trans.\ William Benham, Project Gutenberg eBook 1653, accessed 16 October 2026.} And the Church, keeping vigil for her dead in the Office of the Dead, gives the soul words for its fear. After the seventh lesson of Matins she sings \latin{Peccantem me quotidie, et non me poenitentem, timor mortis conturbat me: Quia in inferno nulla est redemptio, miserere mei, Deus, et salva me}.\endnote{Office of the Dead, Matins, third nocturn, responsory after the seventh lesson, \emph{Rituale Romanum} (Ratisbon: Pustet, 1872), pp.~171--172.} It is the voice of one who sins every day and does not repent, whom the fear of death troubles because in hell there is no redemption, and who therefore cries to God for mercy and salvation while the cry can still be heard.

That is the only use the saints made of eternity: to spend time well. Alphonsus taught his readers to pray, in words he gives as Augustine's, ``Here burn, here cut, here spare not, that you may spare for eternity.''\endnote{Alphonsus, \emph{Preparation for Death}, Consideration XXVII, second point, p.~332.} And Augustine himself, in the midst of his argument for the endlessness of punishment, stopped to draw the only conclusion that matters to a living man: ``they who desire to be rid of eternal punishment ought to abstain from arguing against God, and rather, while yet there is opportunity, obey the divine commands.''\endnote{Augustine, \emph{City of God} XXI.23.} ``Behold, now is the acceptable time: behold, now is the day of salvation.''\endnote{2 Corinthians 6:2, Douay--Rheims (Challoner).}
